\section{Créditos en estudios avanzados: lo que te toca}
El expediente llega en el orden del formato «INTEGRACIÓN DE EXPEDIENTE: CRÉDITOS EN ESTUDIOS AVANZADOS»:
\begin{enumerate}
  \item Dos fotografías tamaño semitítulo, blanco y negro, papel mate, fondo gris claro, de frente, con retoque, varones con camisa blanca, saco y corbata oscuros, mujeres con blazer oscuro
  \item Acta de evaluación profesional
  \item Copia del certificado de bachillerato con firma y sello de cotejo con el original
  \item Constancia de autenticación si el bachillerato no es UAEM \emph{(solo si aplica)}
  \item Dictamen de revalidación expedido por la Secretaría de Educación Pública (sólo si aplica) \emph{(solo si aplica)}
  \item Acta de nacimiento
  \item Certificado total de licenciatura
  \item Solicitud por escrito del egresado dirigida a la Subdirección Académica con firma y sello de la misma
  \item Certificado que ampare el 75 \% de maestría o doctorado integral o el 100 \% del diplomado superior
  \item Dictamen de equivalencia expedido por la Secretaría de Ciencia para cualquier diplomado superior universitario, maestría o doctorado integral
  \item Autorización para continuar con trámites de titulación con fundamento en el artículo 5° del Reglamento de Evaluación Profesional (sólo si aplica) \emph{(solo si aplica)}
  \item Constancia de no adeudo a Servicios Escolares
  \item Formato de revisión para autorización de evaluación profesional
  \item Formato de integración de expediente
  \item Formato de datos del interesado para expedición de título profesional
\end{enumerate}
\subsection{Cómo se genera la carpeta}
El Departamento de Titulación verifica el expediente y lo exporta como carpeta comprimida.
\captura{creditos_servicios_escolares_01_zip}{El Departamento de Titulación exporta la carpeta comprimida que recibe Servicios Escolares.}
